\chapter{Introduction}
\label{chapter1}
\section{Background}
\label{chapter1.Background}

This latex template mainly serves for the MS.c. and Ph.D. student at University of Macau (UM).

This template is reorganized by Yu WAN, a Ph.D. student at NLP$^2$CT lab, University of Macau.


\section{Format of Theses}

The format of thesis at UM follows the instruction file provided by Graduated School (GRS) at UM.\footnote{ \href{https://grs.um.edu.mo/wp-content/uploads/Guidelines-Governing-the-Format-of-PhD-Theses-Final.pdf}{\url{https://grs.um.edu.mo/wp-content/uploads/Guidelines-Governing-the-Format-of-PhD-Theses-Final.pdf}}}

\subsection{Default Format}
The font type thoughout the thesis should be ``Times New Roman'' for English.
The font size should be 12-point (12pt) for almost all cases (including the titles and contents), except the materials aside from main contents (\textit{e.g.}, figures, tables, and footnotes).
The line spacing is double-line spacing for almost all cases except footnotes, bibliographic entries, etc.
The font style should be bolded for most of the titles.

\subsection{Cover, Front, and Colophon Pages}
The materials in cover, font and colophon pages should be centered.

The colophon page should not be revised, and should be single-side when printing.
The beginning of those two paragraphs should not have indents (as the command ``\texttt{noindent}'' used in the \texttt{colophon.tex} file).

The page number should not be presented on those pages.

Also, be patient to revise the related lines in those pages, \textit{e.g.}, your thesis title, your affiliation, etc.
Those contents are marked with the square brackets, and you should modify them \textcolor{red}{based on your own experiences}, \textit{e.g.}, {[YEAR]} $\rightarrow$ 2023.

\subsection{Acknowledgement, Abstract, and Declaration Pages}
The titles of those pages should be centered following the default format.

Except for the Declaration part, the indent should be introduced at the beginning of each paragraph.

The Declaration page should be single-side when printing, \textit{i.e.}, the page number should be an odd integer, and an empty page should be strictly followed.

\subsection{Table of Contents}
The Table of Contents includes the titles and their corresponding page numbers following the default format.

\subsection{List of Tables and Figures}
\label{lotf}
This collects the tables and figures presented in this thesis.

\subsection{Main Text}

The main text contains the several levels of contents.

\subsubsection{Commands}
The commands for managing your manuscript is shown in Table~\ref{table.chapter1.command}.
You should manage your structure of your thesis depending on your own research.
\begin{table*}[t]
    \centering
    \begin{tabular}{ll}
        \toprule
        
        \textbf{Level} & \textbf{Command} \\

        \midrule

        Part & \textbackslash\texttt{part} \\
        Chapter & \textbackslash\texttt{chapter} \\
        Section & \textbackslash\texttt{section} \\
        Subsection & \textbackslash\texttt{subsection} \\
        Subsubsection & \textbackslash\texttt{subsubsection} \\
        Paragraph & \textbackslash\texttt{paragraph} \\

        \bottomrule
    \end{tabular}
    \caption[Command for content levels.]{Command for content levels. For each level, using the corresponding command to declare the titles.}
    \label{table.chapter1.command}
\end{table*}


\subsubsection{Figures and Tables}
You can introduce your figure as in Figure~\ref{fig.chapter1.example_figure}.
Also, all the captions of your figures and tables will be collected into the List of Tables and Figures.
\begin{figure}[t]
    \centering
    \begin{tikzpicture}
        \draw[rectangle, draw=black!30, fill=black!10, very thick, dashed] (0, 0) rectangle node{}  (12.5, 8.5);
        % \node[draw=black!30, fill=black!10, very thick, dashed, align=top] (0, 0) -- (12.5, 0.0) -- (12.5, 8.0) -- (0.0, 8.0) -- (0.0, 0.0) 
        
    \end{tikzpicture}
    \caption[An example figure.]{An example figure. This figure is drawn with tikzpicture package. You can introduce your figure with a given PDF file.}
    \label{fig.chapter1.example_figure}
\end{figure}

If you want to show the short version of your citation in the List of Tables and Figures, comment the following line in the file ``\texttt{preamble.tex}'':
{\texttt{list-entry=heading}}.

Making sure that you have written the short version in your ``{\textbackslash\texttt{caption}}'' command like ``\textbackslash\texttt{caption}{[short caption]}{\{long caption\}}''.

\subsubsection{Footnotes}
Use the command ``\textbackslash\texttt{footnote}'' to manage your footnotes.

\subsubsection{Citation and Reference}
Use the command ``\textbackslash\texttt{cite}'' and ``\textbackslash\texttt{newcite}'' to manage your citations.
For example:

\begin{itemize}
    \item \newcite{vaswani2017attention} presented the \textsc{Transformer} model in the year 2017.
    \item \textsc{Transformer} model \cite{vaswani2017attention} was introduced in the year 2017.
\end{itemize}

The Reference list is shown in the end of thesis following APA format.

\section{Last Comments}
The first part of your thesis should include the following two sections: \uline{Potential Contributions and Statement of Originally}.

Try to receive advices from the stuff in your faculty and revise the format.

If you have any questions, please contact Yu WAN via \texttt{yb87459@um.edu.mo}.

